%% ============================================================
%%  SECCIÓN 9 — Tratamiento de Datos y Resultados
%%  Responsable: Minilux
%% ============================================================

\section{Tratamiento de Datos y Resultados}
\label{sec:tratamiento}

En esta sección se obtienen las magnitudes derivadas a partir de los
datos experimentales de la Sección~\ref{sec:resultados}. Se adopta
$g = \qty{9.81}{\meter\per\second\squared}$.

%% ------------------------------------------------------------------
\subsection{Longitud de contacto del anillo}
%% ------------------------------------------------------------------

La longitud total de contacto es la suma de las circunferencias interna
y externa del anillo:

\begin{equation}
    L = 2\pi\bigl(R_{\text{ext}} + R_{\text{int}}\bigr)
    \label{eq:t9_longitud_contacto}
\end{equation}

Sustituyendo $R_{\text{ext}} = \qty{5.040e-2}{\meter}$ y
$R_{\text{int}} = \qty{4.980e-2}{\meter}$:

\begin{equation*}
    L = 2\pi\,(\num{0.05040} + \num{0.04980}) = \qty{6.2957e-1}{\meter}
\end{equation*}

%% ------------------------------------------------------------------
\subsection{Primer método — masa efectiva y coeficiente de tensión
superficial}
%% ------------------------------------------------------------------

Para cada medición $i$, la masa efectiva equivalente se obtiene
ponderando cada jinetillo por su posición relativa en el brazo mayor:

\begin{equation}
    M_{\text{ef},i} = \sum_{j} m_j \,\frac{n_j}{10}
    \label{eq:t9_masa_efectiva}
\end{equation}

El coeficiente de tensión superficial para cada medición resulta de
combinar la Ec.~\eqref{eq:masa_ef} con la definición
$\gamma = F_{TS}/L$:

\begin{equation}
    \gamma_i = \frac{M_{\text{ef},i}\,g}{L}
    \label{eq:t9_gamma_m1}
\end{equation}

\paragraph{Cálculo representativo — Medición~1}
($J_1 = \qty{20.2}{\gram}$ en $n = 1$;\;
 $J_4 = \qty{1.2}{\gram}$  en $n = 3$):

\begin{equation*}
    M_{\text{ef},1}
    = \num{20.2} \times \frac{1}{10} + \num{1.2} \times \frac{3}{10}
    = \num{2.02} + \num{0.36}
    = \qty{2.38e-3}{\kilogram}
\end{equation*}

\begin{equation*}
    \gamma_1
    = \frac{\qty{2.38e-3}{\kilogram}
            \times \qty{9.81}{\meter\per\second\squared}}
           {\qty{6.2957e-1}{\meter}}
    = \qty{37.09e-3}{\newton\per\meter}
\end{equation*}

Los resultados para las cinco mediciones válidas se consignan en la
Tabla~\ref{tab:t9_gamma_m1}.

\begin{table}[htbp]
\centering
\caption{Masa efectiva y coeficiente de tensión superficial para cada
medición del primer método.}
\label{tab:t9_gamma_m1}
\begin{tabular}{cccc}
\toprule
Medición &
$M_{\text{ef}}$ (\si{\gram}) &
$F_{TS}$ (\si{\milli\newton}) &
$\gamma_i$ (\si{\milli\newton\per\meter}) \\
\midrule
1 & \num{2.38} & \num{23.35} & \num{37.09} \\
2 & \num{2.67} & \num{26.19} & \num{41.60} \\
3 & \num{2.46} & \num{24.13} & \num{38.33} \\
4 & \num{2.58} & \num{25.31} & \num{40.20} \\
5 & \num{2.66} & \num{26.09} & \num{41.45} \\
\bottomrule
\end{tabular}
\end{table}

El valor promedio se obtiene como:

\begin{equation}
    \overline{\gamma}_1 = \frac{1}{N}\sum_{i=1}^{N}\gamma_i
    \label{eq:t9_gamma_promedio}
\end{equation}

\begin{equation*}
    \overline{\gamma}_1
    = \frac{\num{37.09} + \num{41.60} + \num{38.33} + \num{40.20} + \num{41.45}}{5}
    = \qty{39.73}{\milli\newton\per\meter}
\end{equation*}

La propagación de incertidumbre se desarrolla en el
Anexo~\ref{anx:prop_incertidumbre}. El resultado final del primer
método es:

\begin{equation}
    \boxed{\overline{\gamma}_1 = \qty{39.73 \pm 0.88}{\milli\newton\per\meter}}
    \label{eq:t9_resultado_m1}
\end{equation}

El error relativo respecto al valor bibliográfico
$\gamma_{\text{ref}} = \qty{72.8}{\milli\newton\per\meter}$
\cite{serway} es:

\begin{equation*}
    E_r
    = \frac{|\overline{\gamma}_1 - \gamma_{\text{ref}}|}{\gamma_{\text{ref}}}
      \times \qty{100}{\percent}
    = \frac{|\num{39.73} - \num{72.8}|}{\num{72.8}} \times \qty{100}{\percent}
    = \qty{45.4}{\percent}
\end{equation*}

%% ------------------------------------------------------------------
\subsection{Segundo método — coeficiente de tensión superficial}
%% ------------------------------------------------------------------

Aplicando la Ec.~\eqref{eq:gamma_m2} con los datos de la
Tabla~\ref{tab:datos_m2}:
$a = \qty{2.73}{\centi\meter}$,
$b = \qty{2.55}{\centi\meter}$,
$h = \qty{2.78}{\centi\meter}$,
$P = \qty{11.772e-3}{\newton}$.

\paragraph{Cálculo}

\begin{equation*}
    a - b = \num{2.73} - \num{2.55} = \qty{0.18}{\centi\meter},
    \qquad
    a + b = \qty{5.28}{\centi\meter}
\end{equation*}

\begin{equation*}
    \frac{h^2}{a - b} + a + b
    = \frac{(\num{2.78})^2}{\num{0.18}} + \num{5.28}
    = \num{42.94} + \num{5.28}
    = \qty{48.22}{\centi\meter}
    = \qty{4.822e-1}{\meter}
\end{equation*}

Sustituyendo en la Ec.~\eqref{eq:gamma_m2}:

\begin{equation*}
    \gamma_2
    = \frac{P}{2\!\left(\dfrac{h^2}{a-b} + a + b\right)}
    = \frac{\qty{11.772e-3}{\newton}}
           {2 \times \qty{4.822e-1}{\meter}}
    = \frac{\qty{11.772e-3}{\newton}}
           {\qty{9.644e-1}{\meter}}
\end{equation*}

\begin{equation}
    \boxed{\gamma_2 = \qty{12.21}{\milli\newton\per\meter}}
    \label{eq:t9_resultado_m2}
\end{equation}

El error relativo respecto al valor bibliográfico es:

\begin{equation*}
    E_r
    = \frac{|\num{12.21} - \num{72.8}|}{\num{72.8}} \times \qty{100}{\percent}
    = \qty{83.2}{\percent}
\end{equation*}

La marcada diferencia respecto a $\gamma_{\text{ref}}$ es consistente
con la reducción de tensión superficial esperada por efecto del
surfactante, discutida en la Sección~\ref{sec:analisis_exp}.

%% ------------------------------------------------------------------
\subsection{Comparación de resultados}
%% ------------------------------------------------------------------

\begin{table}[htbp]
\centering
\caption{Comparación de los valores de $\gamma$ obtenidos por ambos
métodos con el valor bibliográfico para agua pura a \SI{20}{\celsius}.}
\label{tab:t9_comparacion}
\begin{tabular}{lccc}
\toprule
Método &
$\gamma$ (\si{\milli\newton\per\meter}) &
$E_r$ (\si{\percent}) &
Fluido \\
\midrule
Primer método (Mohr--Westphal) &
    \num{39.73 \pm 0.88} & \num{45.4} & Agua pura \\
Segundo método (película jabonosa) &
    \num{12.21} & \num{83.2} & Solución jabonosa \\
Bibliografía \cite{serway,halliday} &
    \num{72.8} & --- & Agua pura, \SI{20}{\celsius} \\
\bottomrule
\end{tabular}
\end{table}

La Tabla~\ref{tab:t9_comparacion} resume los valores obtenidos por
ambos métodos y su error relativo respecto al valor bibliográfico. El
análisis de las discrepancias observadas se desarrolla en la
Sección~\ref{sec:conclusiones}.